\chapter{实验读数、近似条件与推导依据}
\label{app:checks}

\section{区间计数与谱线宽度}
实验公布的量可能是一格探测器里的计数，也可能是按频率绘出的谱线。与理论比较时，先确定读数对应的区间和单位。同一批事件用动量或波数标记，区间内的计数必须相同，密度的数值则随横坐标改变。算符与矩阵元的记号可查卷首的\hyperref[chap:symbols]{常用符号表}。

\paragraph*{区间概率与概率密度}
对归一化的一维波包，位置探测器第 $j$ 格覆盖 $[x_j,x_j + \Delta x]$。由\secref{sec:continuous}，该格的概率为
\[
    p_j = \int_{x_j}^{x_j+\Delta x}\abs{\psi(x)}^2\,\d x .
\]
$p_j$ 无量纲，因此 $\abs{\psi(x)}^2$ 的量纲是长度的倒数。对同一制备测动量，分别以 $w_p(p)$ 和 $w_k(k)$ 表示每单位动量和每单位波数的概率密度。由 $p = \hbar k$ 及 $w_p(p)\,\d p = w_k(k)\,\d k$，有
\[
    w_k(k) = \hbar\,w_p(\hbar k).
\]
这个因子保证在两种横坐标下，对应区间的面积相同。黑体光谱从频率换到波长也遵循这条规则，具体换算见\ssecref{ss:radiation-density}。

\paragraph*{寿命与半高全宽}
当谱线只受激发态寿命限制时，布居按 $\exp(-\Gamma t)$ 衰减，寿命为 $T_1 = 1/\Gamma$。\exmref{exm:qm-linewidth}得到的强度在角频率 $\omega_0 \pm \Gamma/2$ 处降至峰值的一半。两个半高点之差给出全宽，换到普通频率 $\nu = \omega/(2\pi)$ 得
\begin{equation}\label{eq:qm-lifetime-linewidth}
    \Delta\omega_{\FWHM} = \Gamma,
    \qquad \Delta\nu_{\FWHM} = \frac{1}{2\pi T_1}.
\end{equation}
例如 $T_1 = \qty{10}{ns}$ 对应的寿命展宽约为 $\qty{15.9}{MHz}$。若另有纯退相干，布居寿命可以相同，相干却更快衰减。对\chapref{ch:open}中的指数相干模型，角频率半高全宽为 $2/T_2$，仅当 $T_2 = 2T_1$ 时才等于 $1/T_1$。

\section{物理尺度与近似的适用条件}
对一对孤立能级施加近共振驱动，原子可以多次往返于两态之间。若末态十分密集，同一段驱动却可能表现为近似恒定的跃迁率。选择这两种描述，需要比较耦合强度、能级间隔与观察时间。下面列出几个在正文中实际决定计算方法的尺度。

{\small
\renewcommand{\arraystretch}{1.22}
\begin{longtable}{@{}>{\raggedright\arraybackslash}p{.22\linewidth}@{\hspace{1em}}>{\raggedright\arraybackslash}p{\dimexpr.78\linewidth-1em\relax}@{}}
    \toprule
    问题 & 采用近似的依据及其失效情形\\
    \midrule\endfirsthead
    \toprule 问题 & 采用近似的依据及其失效情形\\\midrule\endhead
    \bottomrule\endfoot
    孤立两态与旋波近似（\secref{sec:rabi}） &
    两态截断要求驱动对其他能级的激发可忽略。对线偏振的近谐振驱动，旋波近似还要求 $\abs{\Omega},\abs{\Delta} \ll \omega$，并在多个载波周期上观察慢演化。正文的理想圆偏振驱动可直接变换到旋转系。\\
    定态微扰（\chapref{ch:approximations}） &
    比较 $\abs{\lambda V_{mn}}$ 与能隙 $\abs{E_n^{(0)} - E_m^{(0)}}$。各一阶混合振幅的模平方之和也应很小。若出现小分母，先把相近能级纳入同一子空间进行对角化。\\
    Fermi 黄金规则（方法~\ref{mtd:qm-golden-window}） &
    设末态间隔为 $\delta E$，矩阵元与态密度组合明显变化的能量尺度为 $W$。要求 $\delta E \ll \hbar/t \ll W$ 且 $\Gamma t \ll 1$。此时观测窗覆盖许多末态，窗内的谱近似平滑，初态的耗尽也可忽略。\\
    自发辐射（\chapref{ch:transitions}） &
    $ka \ll 1$ 允许在尺度为 $a$ 的原子内以同一光场驱动电偶极。自由空间发射率还用到真空辐射模式。放入腔内可能仍满足电偶极条件，却需要重新计算耦合与模式密度。\\
    一阶 Born 散射（\chapref{ch:scattering}） &
    对给定入射能量和出射方向，估计重复散射对振幅的贡献。在共振附近或一阶振幅接近零的方向，高阶项可决定结果。长程势还需单独检查远场边界条件。\\
    绝热跟随（\thmref{thm:qm-adiabatic-following}） &
    对固定的光滑路径，若目标能级全程与其余能级隔开，延长总时长可抑制泄漏。有限时长下，式\eqref{eq:adiabaticcriterion}估计局部混合，还需检验沿整段路径累积的跃迁。\\
    指数弛豫（\chapref{ch:open}） &
    从系统与环境的微观耦合导出无记忆主方程时，常利用弱耦合和环境关联时间远短于系统弛豫时间的条件。若环境在系统弛豫期间仍保留并反馈先前的信息，可把有关的环境自由度纳入模型。\\
    精细结构（\chapref{ch:relativity}） &
    对氢样离子，$Z\alpha \ll 1$ 使相对论修正小于 Coulomb 束缚能。与精密谱线比较还须给定核质量和核尺寸，并考虑所需精度下的辐射修正。\\
\end{longtable}
}

还有两类条件容易与小参数混淆。理想投影测量规定了仪器如何制备测后态，简并子空间内的额外扰动会改变这一更新规则（\ssecref{ss:measurement-record}）。变分法给出的上界对每个满足模型边界条件、能量期望有限的归一化试探态都成立。试探态本身是基态时取等号（\thmref{thm:qm-variational}）。

\section{正文推导与引用的结论}
有些公式在一个具体模型中可以直接推出，推广到一般系统则需要额外的数学结果。下表指出正文的计算在哪一步结束，以及哪一步采用定理或已知解。定理的适用条件随正文链接列出，相关讲义和原始文献见参考资料。

{\small
\renewcommand{\arraystretch}{1.22}
\begin{longtable}{@{}>{\raggedright\arraybackslash}p{.46\linewidth}@{\hspace{1.4em}}>{\raggedright\arraybackslash}p{\dimexpr.54\linewidth-1.4em\relax}@{}}
    \toprule
    正文展开的计算 & 推广时使用的结论\\
    \midrule\endfirsthead
    \toprule 正文展开的计算 & 推广时使用的结论\\\midrule\endhead
    \bottomrule\endfoot
    两个自旋出口的正交投影之和为 $\one$。连续位置测量中，一个像素对应一个区间（\secref{sec:continuous}）。
    & 自伴算符的谱定理（\thmref{thm:qm-spectral-theorem}）保证可以给每个可测读数集合指派谱投影。无界算符还需指定定义域（\secref{sec:selfadjoint}）。\\
    短时幺正演化给出 Schrödinger 方程（式\eqref{eq:schrket}）。时间平移群与自伴生成元间的对应由\thmref{thm:qm-stone}给出，证明见\secref{sec:stone-proof}。
    & 证明使用自伴算符的谱定理，把能量为 $E$ 的各分量乘以 $\exp(-\i Et/\hbar)$。能量与时间演化生成元的物理认定需由实验检验。\\
    从氢的 $1\mathrm{s}$ 态出发，用 $z$ 偏振光激发到 $n = 2$ 能级时，角向积分选出 $2\mathrm{p},m = 0$（表\ref{tab:hydrogen-dipole-z}）。
    & Wigner--Eckart 定理（\thmref{thm:qm-wigner-eckart}）给出一般球张量矩阵元对磁量子数的依赖\cite{littlejohnwe}。\\
    由径向方程的边界条件和级数截断，构造每个允许的 $n,\ell$ 对应的归一化 Coulomb 束缚态（\secref{sec:hydrogen-radial}）。
    & 精细结构计算还引用了 $r^{-2}$ 的期望值，以及 $\ell > 0$ 时 $r^{-3}$ 的期望值\cite{mitfine}。\\
    两个单粒子态相同时，反对称组合的两项相消，由此得到 Pauli 不相容原理（\chapref{ch:identical}）。
    & 哪些粒子应取反对称态由自旋统计定理（\thmref{thm:qm-spin-statistics}）决定。证明需要相对论量子场论的局域性等条件。\\
    在线性扫描的两态模型中，消去一个振幅得到 Weber 方程。比较扫描两端的能量，可确定固定基矢各自对应的能量支（\chapref{ch:adiabatic}）。
    & 从负无穷扫到正无穷的 Landau--Zener 跃迁概率还需要两端渐近解的连接公式\cite{lz}。\\
    分别计算零光子和单光子的探测记录，再对两种结果取平均，得到振幅阻尼信道（\exmref{exm:qm-amplitude-damping}）。
    & 有限维 GKSL 定理（\thmref{thm:qm-gksl}）刻画所有连续、时间齐次、完全正且迹保持的半群生成元\cite{manzano}。\\
\end{longtable}
}

\sectionstar{Stone 定理的证明}
\label{sec:stone-proof}
\thmref{thm:qm-stone}对时间平移群的正则性只要求强连续。证明分两步。先找出可以求导的态，检验它们给出的生成元是否自伴。再用谱定理从生成元构造指数，检验它是否恢复原来的演化。以下用 $\psi$ 简记向量 $\ket{\psi}$，所有向量极限都取 Hilbert 空间的范数。由群律和幺正性，$\norm{\op U(t + \tau)\psi - \op U(t)\psi} = \norm{\op U(\tau)\psi - \psi}$，所以在零时刻的连续性就给出所有时刻的连续性。

\paragraph*{从群构造自伴生成元}
设 $D$ 是式\eqref{eq:qm-stone-generator}中差商 $[\op U(\tau)\psi - \psi]/\tau$ 有极限的全体向量，$\op K\psi$ 表示该极限。差商的线性给出 $D$ 是线性子空间。即使事先不知某个给定的 $\psi$ 能否求导，也可以对它在短时间 $b > 0$ 内的演化取平均，
\[
    \psi_b = \frac1b\int_0^b\op U(s)\psi\,\d s .
\]
积分可看作连续向量的 Riemann 和的极限。由强连续性，$b \to 0$ 时 $\psi_b \to \psi$。由群律，$\op U(\tau)$ 将积分区间平移 $\tau$，两段重叠的部分相消，故
\[
    \frac{\op U(\tau)\psi_b - \psi_b}{\tau}
    = \frac1b\biggp{\frac1\tau\int_b^{b+\tau}\op U(s)\psi\,\d s
       - \frac1\tau\int_0^\tau\op U(s)\psi\,\d s}
    \to \frac{\op U(b)\psi - \psi}{b}.
\]
所以 $\psi_b \in D$，而 $D$ 中的向量可以逼近每个 $\psi$，即 $D$ 稠密。在有限维空间里，每个线性子空间都闭，因此 $D$ 就是全空间。群律还给出
\[
    \dd{t}\op U(t)\phi = \op U(t)\op K\phi = \op K\op U(t)\phi,\qquad \phi \in D ,
\]
因为在任意时刻的差商中，$\op U(t)$ 都可以从左边提出。这也证明 $\op U(t)$ 把 $D$ 映到自身。

对 $\phi,\psi \in D$，幺正性使演化前后的内积相等。在零时刻求导，得到
\[
    0 = \left.\dd{t}\braket{\op U(t)\phi}{\op U(t)\psi}\right|_{t=0}
      = \braket{\op K\phi}{\psi} + \braket{\phi}{\op K\psi}.
\]
按伴随的定义，
$D \subseteq \Dom(\op K\adjoint)$，且在 $D$ 上 $\op K\adjoint = -\op K$。无限维时还需检验两边的定义域相等。
取任意 $\eta \in \Dom(\op K\adjoint)$ 和 $\phi \in D$。把演化移到可微的 $\phi$ 一侧，便能求标量的导数，
\[
    \dd{t}\braket{\phi}{\op U(t)\eta}
    = \dd{t}\braket{\op U(-t)\phi}{\eta}
    = -\braket{\op K\op U(-t)\phi}{\eta}
    = -\braket{\phi}{\op U(t)\op K\adjoint\eta}.
\]
从 $0$ 积到 $t$。等式在与任意 $\phi \in D$ 取内积后都成立，$D$ 又稠密，故得到向量等式
\[
    \op U(t)\eta - \eta = -\int_0^t\op U(s)\op K\adjoint\eta\,\d s.
\]
除以 $t$ 并取 $t \to 0$，右边按强连续性趋于 $-\op K\adjoint\eta$。因此 $\eta \in D$，反向包含也成立，于是
$\op K\adjoint = -\op K$ 包括定义域相等。乘上纯虚常数后，$\op H = \i\hbar\op K$ 就是自伴算符。

\paragraph*{从自伴生成元恢复传播子}
设 $\op H$ 自伴。按\thmref{thm:qm-spectral-theorem}，谱投影 $\op\Pi_H(B)$ 选出能量落在可测集合 $B$ 中的分量。
这个分量的范数平方为 $\mu_\psi(B) = \braket{\psi}{\op\Pi_H(B)\psi}$。特别地，$\mu_\psi(\RR) = \norm{\psi}^2$。谱定理给出
\begin{subequations}\label{eq:qm-propagator-spectral-integral}
    \begin{gather}
        \op V(t) = \int_{\RR}\e^{-\i Et/\hbar}\d\op\Pi_H(E),\\
        \psi \in \Dom(\op H)\  \Longleftrightarrow \
            \int_{\RR} E^2\,\d\mu_\psi(E) < \infty .
    \end{gather}
\end{subequations}
各能量分量只乘模为 $1$ 的相位，相邻两次演化的相位相乘，故 $\op V(t)$ 是幺正群。正交分量的范数平方相加给出
\[
    \norm{\op V(\tau)\psi - \psi}^2
    = \int_{\RR}\abs{\e^{-\i E\tau/\hbar} - 1}^2\,\d\mu_\psi(E).
\]
在任何有限能量区间 $[-M,M]$ 上，被积函数一致趋于零。区间外的积分至多为 $4\mu_\psi(\abs E > M)$，取 $M$ 充分大即可任意小。因此 $\op V(t)$ 强连续。

对差商做同样的估计。每个能量分量上的差商及其极限为
\[
    q_\tau(E) = \frac{\e^{-\i E\tau/\hbar} - 1}{\tau}
        \ \xrightarrow{\tau\to0}\  - \frac{\i E}{\hbar}.
\]
在每个有限能量区间上，这个收敛是一致的。区间外，由 $\abs{\e^{\i x} - 1} \leq \abs x$ 得 $\abs{q_\tau(E)} \leq \abs E/\hbar$，故误差满足
\[
    \int_{\abs E>M}\biggabs{q_\tau(E) + \frac{\i E}{\hbar}}^2\d\mu_\psi(E)
        \leq \frac{4}{\hbar^2}\int_{\abs E>M}E^2\,\d\mu_\psi(E).
\]
当 $\psi \in \Dom(\op H)$ 时，取 $M$ 充分大可使右侧任意小。因此 $\op V(t)\psi$ 在零时刻的导数为 $-\i\op H\psi/\hbar$。
反之，若 $\op V(t)\psi$ 在零时刻有导数 $\xi$，将差商投影到 $[-M,M]$，在这个区间上求极限得
\[
    \int_{[-M,M]} E^2\,\d\mu_\psi(E)
    = \hbar^2\norm{\op\Pi_H([-M,M])\xi}^2 \leq \hbar^2\norm{\xi}^2 .
\]
右侧与 $M$ 无关，所以全能谱上的二阶矩有限，$\psi$ 必属于 $\Dom(\op H)$。

最后，把 $\op V(t)$ 用于第一步从 $\op U(t)$ 得到的 $\op H$。两种演化保持同一定义域，其上的每条轨道都服从式\eqref{eq:schrket}。对 $\phi,\psi \in \Dom(\op H)$，内积
$\braket{\op U(t)\phi}{\op V(t)\psi}$ 的导数为零，因为两项中的 $\op H$ 可由自伴性移到同一侧，而系数互为相反数。
这个内积因此等于初值 $\braket{\phi}{\psi}$。由定义域稠密及幺正性，等式延伸到任意两向量，于是 $\op U(t)\adjoint\op V(t) = \one$，即 $\op U(t) = \op V(t)$。而给定 $\op U(t)$，其零时刻的差商已经唯一确定 $\op H$ 的定义域和作用，定理得证。

\section{计算结果的独立检验}
检验一个近似结果，可以选用另一种可计算的极限，或选取在推导中未曾用过的守恒量。下面三个例子分别针对演化时间、散射通量和能谱截断，它们检验的是不同类型的误差。

\paragraph*{用短时极限核对驱动强度}
在\secref{sec:rabi}的两态模型中，把式\eqref{eq:rabi}在 $\Omega_{\mathrm{R}} t \ll 1$ 处展开，得到
\[
    P_{-z\to+z}(t) = \frac{\Omega^2t^2}{4} + \bigo(t^4).
\]
独立地，驱动的非对角矩阵元大小为 $\hbar\abs{\Omega}/2$，故一阶跃迁振幅的模在短时为 $\abs{\Omega}t/2$。两种算法给出相同的二次项，可检出驱动幅度定义中相差的因子 $2$。

\paragraph*{用通量核对透射概率}
在\secref{sec:tunneling}的一维定态散射中，反射和透射概率均是出射通量与入射通量之比。取无吸收的实势，仅左右两个传播通道开放。记入射和透射速度分别为 $v_{\mathrm{L}},v_{\mathrm{R}} > 0$，透射振幅与入射振幅之比为 $t$，则
\[
    T = \frac{v_{\mathrm{R}}}{v_{\mathrm{L}}}\abs t^2,\qquad R + T = 1 .
\]
两侧势能不同时，遗漏速度比会使算出的总概率偏离 $1$。

\paragraph*{用完整算符核对截断本征态}
记保留的基矢所张成的子空间为 $S$，到 $S$ 的正交投影为 $\op\Pi$。在 $S$ 上对角化 $\op\Pi\op H\op\Pi$，得到归一化本征态 $\ket{\phi}$ 和本征值 $E$。若 $\ket{\phi}$ 属于 $\op H$ 的定义域，它在完整空间中的残差为
\[
    (\op H - E\one)\ket{\phi} = (\one - \op\Pi)\op H\ket{\phi}.
\]
左式和右式相等，因为保留空间内的分量已在对角化时抵消。残差范数给出 $E$ 到完整能谱距离的上界。若要判定附近的谱值属于哪个能级，还需要相邻能级的间隔（\dfnref{dfn:qm-eigen-residual}）。

\section{典型习题与历年试题}
\label{sec:qm-exam-problems}
题名链接到本卷的题干和解答。右栏的公开试题考查相同的物理问题或解题方法，具体模型、参数和设问可有差异。

{\small
\renewcommand{\arraystretch}{1.2}
\begin{longtable}{@{}r@{\hspace{.8em}}>{\raggedright\arraybackslash}p{.40\linewidth}@{\hspace{1em}}>{\raggedright\arraybackslash}p{.49\linewidth}@{}}
    \toprule
    页 & 本卷例题 & 公开试题\\
    \midrule\endfirsthead
    \toprule 页 & 本卷例题 & 公开试题\\\midrule\endhead
    \bottomrule\endfoot
    \pageref{exm:qm-exam-sudden-well} & \nameref{exm:qm-exam-sudden-well} &
    \href{https://gs.njust.edu.cn/_upload/article/files/0c/cb/e2112c3842e1a5335223fc2c548c/58c5ce66-f97f-4e69-b19b-0121747d961f.pdf}{南京理工大学，2017 年硕士入学考试 843，第六题}\\
    \pageref{exm:qm-exam-wronskian} & \nameref{exm:qm-exam-wronskian} &
    \href{https://www.isibang.ac.in/~statmath/oldqp/MQP/Quantum\%20Mechanics\%20I\%20MS\%202024-25.pdf}{印度统计研究所，2025 年量子力学期中，第 1(a) 题}\\
    \pageref{exm:qm-exam-half-oscillator} & \nameref{exm:qm-exam-half-oscillator} &
    \href{https://ocw.mit.edu/courses/8-04-quantum-physics-i-spring-2013/edece27b26b6c7d23f6657eeaa3b09b7_MIT8_04S13_finalexam.pdf}{MIT，2013 年春 8.04 期末，第 6(b) 题}\\
    \pageref{exm:qm-cartesian-angular-readout} & \nameref{exm:qm-cartesian-angular-readout} &
    \href{https://physics.unl.edu/sites/unl.edu.cas.physics/files/media/file/DayII_Aug2022.pdf}{内布拉斯加大学林肯分校，2022 年 8 月预考，量子力学 B1}\\
    \pageref{exm:qm-degenerate-dark-combination} & \nameref{exm:qm-degenerate-dark-combination} &
    \href{https://gs.njust.edu.cn/_upload/article/files/f1/73/eca4fb0a437fac976c1330a5189f/37a3616e-61d5-4ace-81a2-ac2fe40dcf84.pdf}{南京理工大学，2022 年硕士入学考试 843，计算题第 1 题}\\
    \pageref{exm:qm-exam-weak-binding} & \nameref{exm:qm-exam-weak-binding} &
    \href{https://phys.washington.edu/sites/phys/files/documents/grad/phd_mreprobs.pdf}{华盛顿大学，2013/14 学年硕士考试，量子力学第 1B 题}\\
    \pageref{exm:qm-curvature-pulse} & \nameref{exm:qm-curvature-pulse} &
    \href{https://phys.washington.edu/sites/phys/files/documents/grad/phd_mreprobs.pdf}{华盛顿大学，2018 年硕士考试，量子力学第 4 题}\\
    \pageref{exm:qm-fermion-contact-box} & \nameref{exm:qm-fermion-contact-box} &
    \href{https://www.physics.uci.edu/sites/default/files/Spring_2017_1.pdf}{加州大学尔湾分校，2017 年春博士资格考，量子力学第 1 题}\\
    \pageref{exm:qm-exam-two-centres} & \nameref{exm:qm-exam-two-centres} &
    \href{https://pa.msu.edu/sites/_pa/assets/File/QMS18final_no\%20solutions.pdf}{密歇根州立大学，2018 年春量子力学科目考试，第 3 题}\\
\end{longtable}
}
